\subsection*{三种次序从同一状态开始}

以下三条历史都从第九章已发布 EV1 的相同状态开始：A 已保存 P1/C1/T1，平台为 Q7 的 CNY 3000 分建立 X1，付款账户由 10000 分变为 7000 分。R 已建立 O1、保存原结果与 EV1，P1 为 \code{paid}，T1 为 \code{done}。原购买身份和快照保持，EV1 已进入指定的 \code{fulfillment.paid} 队列，履约库没有 F，也没有消费记录。

这里的 ACK 指消费者向消息代理发送的处理确认。手工确认模式下，代理收到确认以后，可以把对应投递标为可删除。何时具备确认依据，由消费应用决定~\cite{rabbitmq-reliability,rabbitmq-confirms}。第九章的 publisher confirm 则是代理给发布者的受理回报。发布受理与本次消费各有自己的进度。

先确认，再处理。消费者收到 EV1，立即发送 ACK，代理保存了该进度。消费者随后准备插入 F，尚未提交就退出。恢复后履约库为空，代理已经可以移除原副本。这份消息的恢复责任已经在业务提交以前结束，数据库和消费者恢复可用后，也可能再没有输入触发建单。

再加一张已处理事件表。消费者先独立提交 \code{EV1 已处理}，然后才准备建立 F。它在两次提交之间退出。代理重投 EV1，恢复者查到已处理记录，照常跳过并确认。数据库最终保存着一份已处理记录，O1 的初次履约却停在建单以前。

最后把顺序反过来：先提交 F1，再单独记录 EV1。消费者在记录以前退出。恢复者没有找到消费记录，又插入 F2。在这个尚无业务唯一约束的错误方案中，同一 O1 得到两张初次履约单。若后续调度各执行一次，重复建单还可能继续向外传播。

\CTenFigure{ch10-original01}{三条共用初态的退出历史：先确认、先存记录与先建履约单}{fig:ch10-three-orders}

三个问题分别来自两项事实被拆开提交。先 ACK 时，代理进度领先于本地效果。先单独存记录时，处理标记领先于效果。先单独修改业务时，效果领先于记录。要让后来者依据记录恢复，记录与它所证明的本地效果就需要共同提交。

\subsection*{投递标记说明了哪个对象}

RabbitMQ 为一次投递分配 \code{delivery tag}，它在同一 channel 内标识投递，确认也必须发回原 channel~\cite{rabbitmq-confirms}。重连建立新 channel，编号可以重新开始\footnote{RabbitMQ 文档用短语 ``scoped per channel'' 说明投递编号的范围。连接恢复建立新 channel 后，消费者按新 channel 的投递与编号发送确认。}。因此消费者把 channel 与 tag 一起用于确认当前投递，把消息中的稳定事件身份用于恢复业务结果。

\code{redelivered} 提示某个代理消息副本曾经投递后重新入队。它为 \code{true} 时，当前消费者可能第一次见到该副本，因为先前接收者可以是另一实例。它为 \code{false} 时，说明这个副本尚未经历这种重投~\cite{rabbitmq-reliability}。第九章的 PA1 和 PA2 却可以分别产生 M1、M2，两份都携带 EV1。M1 已经建立 F1，M2 第一次交付时可以是 \code{redelivered=false}。

因此每个到达的合法业务消息都进入同一消费判断。若只在 \code{redelivered=true} 时查记录，M2 会绕开防重。发布尝试 PA、代理副本、投递 tag 与业务事件 EV 所标识的对象，需要沿这条具体历史逐层区分。
